\documentclass[10pt,a4paper]{article}
\usepackage{amsmath,amssymb,amsthm,graphicx,hyperref,xspace,listings,mathtools,tikz-cd,csquotes,etoolbox,xparse}
\usepackage[capitalize]{cleveref}
\usetikzlibrary{calc,arrows.meta,decorations.pathreplacing,calligraphy,backgrounds}
\newtheorem{theorem}{Theorem}
\newtheorem{lemma}[theorem]{Lemma}
\newtheorem{proposition}[theorem]{Proposition}
\newtheorem{definition}[theorem]{Definition}
\newtheorem{remark}[theorem]{Remark}
\newtheorem{example}[theorem]{Example}
\sloppy
%%% couleurs
\definecolor{egreen}{rgb}{0, 0.4, 0.267}
\definecolor{dkviolet}{rgb}{0.6,0,0.8}
\definecolor{dkgreen}{rgb}{0,0.4,0}
\definecolor{dkblue}{rgb}{0,0.1,0.5}
\definecolor{lightblue}{rgb}{0,0.5,0.5}
\definecolor{orange}{rgb}{0.9,0.39,0}
\definecolor{lightgrey}{RGB}{240,240,240}

% Names of systems/libraries
\newcommand{\namefont}[1]{\textsf{#1}}
\newcommand{\Coq}{\namefont{Coq}\xspace}
\newcommand{\Lean}{\namefont{Lean}\xspace}
\newcommand{\Globular}{\namefont{Globular}\xspace}
\newcommand{\MC}{\namefont{Mathematical Components}\xspace}

\newcommand{\FANL}{\namefont{FANL}\xspace}

% Notes
\newcommand{\anote}[1]{\todo[color=blue!50]{Assia: #1}}
\newcommand{\mnote}[1]{\todo[color=green!50]{{\scriptsize Matthieu}: #1}}
\newcommand{\anoteinl}[1]{\todo[color=blue!50, inline]{Assia: #1}}
\newcommand{\mnoteinl}[1]{\todo[color=green!50, inline]{Matthieu: #1}}

%Latin abbrev.
\newcommand{\ie}{i.e.}
\newcommand{\eg}{e.g.}
\newcommand{\resp}{resp.\ }

%Various abbrev symbol commands
\renewcommand{\~}{\widetilde}
\renewcommand{\emptyset}{\varnothing}

\newcommand{\simto}{\xrightarrow{\raisebox{-3pt}[0pt][0pt]{\small$\hspace{-1pt}\sim$}}}

% Numbers
\newcommand{\Z}{\mathbb Z}
\newcommand{\N}{\mathbb N}

% Mathcal
\newcommand{\set}{\mathcal S}
\newcommand{\Cat}{\mathcal C}
\newcommand{\G}{\mathcal G}
\newcommand{\T}{\mathcal T}
\newcommand{\M}{\mathcal M}
\newcommand{\A}{\mathcal A}
\newcommand{\Q}{\mathcal Q}
\newcommand{\BP}{\mathcal{BP}}
\newcommand{\BT}{\mathcal{BT}}
\newcommand{\E}{\mathcal E}

%Mathrm
\newcommand{\x}{\mathrm{x}}
\newcommand{\PQ}{\mathrm{PQ}}
\newcommand{\im}{\mathrm{Im}}
\newcommand{\card}[1]{\mathrm{card}(#1)}
\newcommand{\Ob}{\mathrm{Ob}}
\newcommand{\Hom}{\mathrm{Hom}}
\newcommand{\restr}{\mathrm{restr}}
\newcommand{\com}{\mathrm{commute}}
\newcommand{\id}{\mathrm{id}}

\newcommand{\cone}{\mathrm{cone}}
\newcommand{\comp}{\mathrm{comp}}
\newcommand{\bi}{\mathrm{bi}}
\newcommand{\map}{\mathrm{map}}
\newcommand{\qdot}{\mathrm{dot}}
\newcommand{\cat}{\mathrm{cat}}
\newcommand{\ab}{\mathrm{ab}}

%Mathsf
\newcommand{\sop}{\textsf{sp}}
\newcommand{\tp}{\textsf{tp}}
\newcommand{\sot}{\textsf{st}}
\newcommand{\tot}{\mathrm{tot}}

%Op
\newcommand{\free}[1]{{\langle #1 \rangle}}
% \newcommand{\rel}[1]{\overset{#1}{\sim}}
% \newcommand{\rel}[1]{\sim_{#1}}
\newcommand{\rel}[1]{\sim}
\newcommand{\eqd}{\approx}
\newcommand{\cyc}{\mathring}
\newcommand{\dual}{\dagger}

\newcommand{\RestrComp}{\mathrm{RestrComp}}
\newcommand{\RestrId}{\mathrm{RestrId}}
\newcommand{\Pushout}{\mathrm{Cospan}}
\newcommand{\PushoutE}{\mathrm{PushoutEU}}
\newcommand{\Comp}{\mathrm{Comp}}
\newcommand{\CompE}{\mathrm{CompE}}
\newcommand{\ComEq}{\mathrm{ComEq}}
\newcommand{\EqId}{\mathrm{EqId}}
\newcommand{\EqPath}{\mathrm{EqPath}}
\newcommand{\EqPathRefl}{\mathrm{EqPathRefl}}
\newcommand{\EqPathSym}{\mathrm{EqPathSym}}
\newcommand{\EqPathTrans}{\mathrm{EqPathTrans}}
\newcommand{\PathCom}{\mathrm{PathCom}}
\newcommand{\EqPathConcat}{\mathrm{EqPathConcat}}
\newcommand{\Id}{\mathrm{Id}}
\newcommand{\IdE}{\mathrm{IdE}}
\newcommand{\EmptyEU}{\mathrm{EmptyEU}}
\newcommand{\CommergeWithId}{\mathrm{CommergeWithId}}

\newcommand{\cl}{c\ell} % path-relation closure
\newcommand{\concat}{\!\cdot\!} % concatenation of paths
\newcommand{\acc}{\mathrm{acc}} % acc(u) ::= vertices accessible from u
\newcommand{\arcsfrom}[1]{\A_{#1}} % acc(u) ::= adjacent edges to u going to v
\newcommand{\arcsto}[1]{\;\overline{\!\!\A}_{#1}} % acc(u) ::= adjacent edges to v from u


%Mathrm formula names
\newcommand{\Mono}{\mathrm{Mono}}
\newcommand{\Epi}{\mathrm{Epi}}
\newcommand{\Cone}{\mathrm{Cone}}
\newcommand{\Limit}{\mathrm{Limit}}
\newcommand{\Colimit}{\mathrm{Colimit}}
\newcommand{\Zero}{\mathrm{Zero}}
\newcommand{\ZeroE}{\mathrm{ZeroE}}
\newcommand{\Ker}{\mathrm{Ker}}
\newcommand{\Coker}{\mathrm{Coker}}
\newcommand{\ProductE}{\mathrm{ProductE}}
\newcommand{\CoproductE}{\mathrm{CoproductE}}
\newcommand{\KerE}{\mathrm{KerE}}
\newcommand{\CokerE}{\mathrm{CokerE}}
\newcommand{\MonoNormal}{\mathrm{MonoNormal}}
\newcommand{\EpiNormal}{\mathrm{EpiNormal}}
\newcommand{\Commerge}{\mathrm{Commerge}}


% Drawing quivers, also to be used in formulas

\newcommand{\rawvertex}[1]{\tikz{\fill[#1] (0,0) circle (2pt);}}
\newcommand{\vertex}[1]{node {\rawvertex{#1}}}

% See \cobimapimono for a detailed example
% No first option (\ifblank is false): subscript form
% Empty first option (\ifblank is true): normal form
\NewDocumentCommand{\quiver}{ O{1} m O{0} m O{0} }{
  \tikz[baseline={([yshift=\ifblank{#1}{-.6ex}{0ex}]current bounding box.center)}, inner sep=\ifblank{#1}{1pt}{.5pt}, -{Latex[scale=\ifblank{#1}{1}{.5}]}, scale=\ifblank{#1}{1}{.5}]{
    \bfseries
    \foreach \p [count = \j, evaluate={\i=int(\j-1)}, evaluate={\k=int(100-(\i<#3)*70)}] in {#2} {
      \node[black!\k] (\i) at \p {.};
    }
    \foreach \i/\j [count = \l, evaluate={\k=int(100-(\l<=#5)*70)}] in {#4} {
      \draw[black!\k] (\i)--(\j);
    }
  }
}

\NewDocumentCommand{\quiverr}{ O{1} m O{0} m O{0} }{
  \tikz[baseline={([yshift=\ifblank{#1}{-.5ex}{0ex}]current bounding box.center)}, inner sep=\ifblank{#1}{1pt}{.5pt}, -{Latex[scale=\ifblank{#1}{1}{.5}]}, scale=\ifblank{#1}{1}{.5}]{
    \bfseries
    \foreach \p [count = \j, evaluate={\i=int(\j-1)}, evaluate={\k=int(100-(\i<#3)*70)}] in {#2} {
      \node[black!\k] (\i) at \p {.};
    }
    \foreach \i/\j/\s [count = \l, evaluate={\k=int(100-(\l<=#5)*70)}] in {#4} {
      \draw[black!\k] (\i)to[bend left=\s](\j);
    }
  }
}


\newcommand{\Qdot}[1][1]{{\quiver[#1]{(0,0)}{}}}\newcommand{\Qtwodots}[1][1]{{\quiver[#1]{(0,0),(.2,0)}{}}}
\newcommand{\Qloop}{\tikz[baseline=0pt, inner sep=.5pt, scale=.5]{\node (0) at (0,0) {\textbf{.}}; \draw[-{Latex[scale=.5]}] (0) to[out=120, in=60, looseness=10] (0);}}
\newcommand{\Qmap}[1][1]{{\quiver[#1]{(0,0),(.8,0)}{0/1}}}
\newcommand{\Qbimap}[1][1]{{\quiver[#1]{(0,0),(1,0),(2,0)}{0/1,1/2}}}
\newcommand{\Qcobimap}[1][1]{{\quiverr[#1]{(0,0),(1,0)}{0/1/30,0/1/-30}}}
\newcommand{\Qcomp}[1][1]{{\quiver[#1]{(0,0),(.5,-.5),(1,0)}{0/1,0/2,1/2}}}
\newcommand{\Qmono}[1][1]{{\quiverr[#1]{(0,0),(.8,0),(1.6,0)}{0/1/20,0/1/-20,0/2/35,1/2/0}}}
\newcommand{\Qker}[1][1]{{\quiver[#1]{(0,0),(.7,.2),(.9,-.2),(1.6,0)}{0/1,0/2,0/3,1/3,2/3}}}

% naming convention : source quiver, number in binary from i, target quiver ; the number k indicates the k-th such embedding in lexicographical order
\newcommand{\mapicomp}{{\quiver{(1,0),(.5,-.5),(0,0)}[1]{1/0,2/0,2/1}[2]}}
\newcommand{\mapiocomp}{{\quiver{(.5,-.5),(0,-0),(1,0)}[1]{1/0,0/2,1/2}[2]}}
\newcommand{\mapiicomp}{{\quiver{(0,0),(.5,-.5),(1,0)}[1]{0/1,0/2,1/2}[2]}}
\newcommand{\dotimap}{{\quiver{(1,0),(0,0)}[1]{1/0}[1]}}
\newcommand{\dotiomap}{{\quiver{(0,0),(1,0)}[1]{0/1}[1]}}
\newcommand{\twodotsimap}[1][1]{{\quiver[#1]{(0,0),(1,0)}{0/1}[1]}}
\newcommand{\mapioomono}  {{\quiverr{(0,0),(.8,0),(1.6,0)}[1]{0/1/20,0/1/-20,0/2/35,1/2/0}[3]}}
\newcommand{\compimono}   {{\quiverr{(0,0),(.8,0),(1.6,0)}{0/1/-20,0/1/20,0/2/35,1/2/0}[1]}}
\newcommand{\compiomono}  {{\quiverr{(0,0),(.8,0),(1.6,0)}{0/1/20,0/1/-20,0/2/35,1/2/0}[1]}}
% \cobimapimono has three vertices positioned as in first argument, it can only be used in subscript mode. See the definitions of \Qsomething above to see the way to use it in both forms.
% It has four arrows, the first one for instance goes from the 2-nd vertex to the 0-th and the arrow is bend to the left by an angle of 35 degrees.
% The first option impose the 0-th vertex to be gray (the other are black), the second option impose the two first arrows to be gray. No option is equivalent to [0].
\newcommand{\cobimapimono}{{\quiverr{(1.6,0),(.8,0),(0,0)}[1]{2/0/35,1/0/0,2/1/20,2/1/-20}[2]}}
\newcommand{\dotiiker}{{\quiver{(0,0),(.7,.2),(1.6,0),(.9,-.2)}[3]{0/1,0/3,0/2,1/2,3/2}[5]}}
\newcommand{\mapiooker}{{\quiver{(0,0),(.9,-.2),(.7,.2),(1.6,0)}[2]{0/2,0/1,0/3,1/3,2/3}[4]}}
\newcommand{\mapiker}{{\quiver{(.9,-.2),(1.6,0),(0,0),(.7,.2)}[2]{2/0,2/1,3/1,0/1,2/3}[4]}}
\newcommand{\veeiker}{{\quiver{(0,0),(.7,.2),(.9,-.2),(1.6,0)}[1]{0/1,0/2,0/3,1/3,2/3}[3]}}
\newcommand{\mapibimap}{\quiver{(2,0),(1,0),(0,0)}[1]{1/0,2/1}[1]}
\newcommand{\mapiobimap}{\quiver{(0,0),(1,0),(2,0)}[1]{0/1,1/2}[1]}


% Embeddings with paths
\newcommand{\linebox}{\path[use as bounding box] (-.05,0ex)--(.55,0ex)}
\newcommand{\spo} {\tikz[scale=.5, thick]{\linebox; \draw[black!30] (0,0) node[black] {\textbf{.}} to[out=40,in=-140] (.5,0);}}
\newcommand{\tpo} {\tikz[scale=.5, thick]{\linebox; \draw[black!30] (0,0) to[out=40,in=-140] (.5,0) node[black] {\textbf{.}};}}
\newcommand{\spd} {\tikz[scale=.5, thick]{\linebox; \draw[black!30] (0,0) to[out=40,in=-140] (.5,0); \clip (0,0) circle (.23);  \draw (0,0) to[out=40,in=-140] (.5,0);}}
\newcommand{\tpd} {\tikz[scale=.5, thick]{\linebox; \draw[black!30] (0,0) to[out=40,in=-140] (.5,0); \clip (.5,0) circle (.23); \draw (0,0) to[out=40,in=-140] (.5,0);}}
\newcommand{\tpdi}{\tikz[scale=.5, thick]{\linebox; \draw[black!30] (0,0) to[out=40,in=-140] (.5,0); \clip (.5,0) circle (.3);  \draw (0,0) to[out=40,in=-140] (.5,0);}}
\newcommand{\std} {\tikz[scale=.5, thick]{\linebox; \draw[black!30] (0,0) node[black] {\textbf{.}} to[out=40, in=-140] (.5,0) node[black] {\textbf{.}};}}


% Inlined code highlighting, wih or without LaTeX escaping
\newcommand{\coq}[1]{\lstinline[mathescape=false]`#1`}
\newcommand{\coqescape}[1]{\lstinline[mathescape=true,escapeinside={/}{/}]`#1`}

% Native minted listing float retained without an external highlighter.
\usepackage{float}
\newfloat{listing}{tbp}{lol}
\floatname{listing}{Listing}
